%! Characteristics of Quadratic Functions — Unit 2, Lesson 2.0 — Warm-Up
% 12pt like every other student component, and it must fit EXACTLY ONE PAGE, blank and keyed.
% Three spiral-review items rehearsing exactly what today leans on:
%   1. substitution with negatives and squares — and it quietly produces the anchor curve's
%      zeros (-1, 3) and its vertex output (-4) before the graph is ever shown;
%   2. reading a LINEAR graph's features (Unit 1), so "read the graph, name the feature" is
%      already the move;
%   3. the hand-off, LEFT HANGING: a pre-drawn parabola with its high point circled — is the
%      "highest" answer the x-value, the y-value, or both? Notes row 1 answers the first two
%      blanks; row 4 settles the last one.
% Both figures are PRE-DRAWN and read. Nothing here asks a student to sketch a graph.
% PAGE LOCKSTEP with warmup_key/: every \blank{} becomes an \ans{}; nothing else differs.
\documentclass[12pt]{article}
\usepackage{algebra2-article}
\usepackage{algebra2-boxes}
% Coordinate grid: {xmin}{xmax}{ymin}{ymax}
\newcommand{\lgrid}[4]{%
  \draw[step=1, linegray!40, very thin] (#1,#3) grid (#2,#4);
  \draw[->, semithick, charcoal] (#1,0)--(#2,0) node[below right, font=\scriptsize]{$x$};
  \draw[->, semithick, charcoal] (0,#3)--(0,#4) node[left, font=\scriptsize]{$y$};}
\newcommand{\fdot}[3]{\fill[#1] (#2,#3) circle (3.2pt);}

\begin{document}

\pageheader{Unit 2, Lesson 2.0}{Warm-Up}
% NAMESTRIP: no \namedateperiod here — the student writes their name once, on the cover.

\vspace{0.10in}

\boxguard[12]
\begin{notesbox}{Getting Ready: Skills We'll Use Today}
Three things you already know. Today we meet a graph that \emph{turns around}, and these are the
tools you'll need --- work quickly, then check with a neighbour.

\vspace{6pt}
\textbf{1. Evaluate the rule.} For $f(x)=x^2-2x-3$, find each output. Watch the negatives.
\begin{center}
\renewcommand{\arraystretch}{1.3}
\begin{tabular}{|c|c|c|c|c|}
\hline
$x$ & $-1$ & $0$ & $1$ & $3$ \\ \hline
$f(x)=x^2-2x-3$ & \blank{1.2cm} & \blank{1.2cm} & \blank{1.2cm} & \blank{1.2cm} \\ \hline
\end{tabular}
\end{center}
Which two inputs gave an output of $0$? \blank{2.4cm}

\vspace{6pt}
\textbf{2. Read the line (Lesson 1.1).} The graph of $y=-2x+4$ is drawn for you.
\begin{minipage}[c]{0.60\linewidth}
  $x$-intercept: \blank{1.6cm} \qquad $y$-intercept: \blank{1.6cm} \\[5pt]
  Is the line increasing or decreasing? \blank{2.6cm} \\[5pt]
  Domain: \blank{3.6cm}
\end{minipage}\hfill
\begin{minipage}[c]{0.36\linewidth}\centering
\begin{tikzpicture}[scale=0.40]
  \lgrid{-1.5}{3.5}{-2.5}{5.5}
  \draw[very thick, forest] (-0.5,5)--(3,-2);
  \fdot{forest}{2}{0}
  \fdot{forest}{0}{4}
\end{tikzpicture}
\end{minipage}

\vspace{6pt}
\textbf{3. A graph that turns around.} This one is \emph{not} a line. Its highest point is
circled.
\begin{minipage}[c]{0.60\linewidth}
  It is increasing until $x=$ \blank{1.2cm}, and decreasing after $x=$ \blank{1.2cm}. \\[5pt]
  Its highest point is at \blank{2.0cm} \\[5pt]
  \textbf{Is the ``highest'' answer the $x$-value, the $y$-value, or both?} \\[3pt]
  \blank{6.4cm}
\end{minipage}\hfill
\begin{minipage}[c]{0.36\linewidth}\centering
\begin{tikzpicture}[scale=0.40]
  \lgrid{-1.5}{5.5}{-2.5}{5.5}
  \draw[very thick, forest, domain=-0.4:4.4, samples=60, smooth]
    plot (\x, {-(\x-2)*(\x-2)+4});
  \fdot{forest}{2}{4}
  \draw[very thick, redacc] (2,4) circle (7pt);
\end{tikzpicture}
\end{minipage}
\end{notesbox}

\end{document}
